\partdivider{B}{The Laws of Motion}{Kinematics describes motion; this part explains it. One law, $\vF=m\va$, applied with a disciplined free-body diagram, accounts for pulleys, friction, banked roads, loops and lifts.}{\chref{ch:newton} Newton's Laws\\[1.2mm]
\chref{ch:connected} Connected Bodies\\[1.2mm]
\chref{ch:friction} Friction\\[1.2mm]
\chref{ch:circdyn} Circular Dynamics and Non-Inertial Frames}

\ncertchapter{7}{Newton's Laws and Free-Body Diagrams}\chlabel{ch:newton}

Part A described motion. This chapter begins to explain it, and the explanation fits in three
sentences published by Newton in 1687, in the \emph{Principia}. The first law says that a body left alone keeps moving uniformly,
which defines what it means for a force to be absent. The second says that a force changes a body's
momentum (its mass times its velocity) at a rate equal to the force, which in most problems means
$\vF=m\va$. The third says that forces come in equal and opposite pairs acting on two different bodies.
Everything in the rest of this Part, and most of what follows it, applies those three sentences.

Of the three, the second is the one that produces answers, and only when it is applied with a particular
discipline: isolate one body, draw every real force acting on it and nothing else, resolve, and only
then write the equation. That discipline, the free-body diagram, is the real subject of the chapter.
Along the way come impulse and momentum conservation, which follow directly from the second and third
laws and will carry \chref{ch:collisions}, and the three commonest misconceptions: that a moving body
needs a force to keep moving, that an action and its reaction cancel, and that a scale measures weight.

\section{The Three Laws}

Each law is stated in a modern paraphrase of Newton's wording and then examined for what it actually
asserts. The words the laws rely on (force, net force, inertial frame, mass, momentum, weight) are
defined as they are needed.

\subsection{The first law: inertia, and what force \emph{is}}

For nearly two thousand years the dominant view was Aristotle's: a moving body stops unless something
keeps pushing it. Push a book across a table and stop pushing; it stops, and Aristotle seems right. The
claim that bodies continue for ever unless something interferes contradicts everyday experience, and
becomes believable only once one notices an unseen agent, friction, present every time one checks. Galileo argued for it from experiments with balls rolling on inclined
planes, Descartes stated it for motion in a straight line, and Newton made it the first of his laws.

\begin{keybox}{Newton's first law (law of inertia)}
Every body continues in its state of rest or of uniform motion in a straight line unless compelled by a
net external force to change that state.
\end{keybox}

The first law does two subtle jobs. It \emph{defines force} qualitatively: a \textbf{force} is a push
or pull exerted on a body by some other body (\emph{external} means exerted by something outside the
body), and it is what changes the body's state of motion. Some forces act through contact (a hand, a
surface, a string); gravity acts across empty space. Force is a vector, written $\vF$ and measured in
newtons (defined below), and several forces on one body add by the triangle law of \chref{ch:vectors};
only their resultant, the \textbf{net force} $\vF_{\text{net}}$, affects the motion.

The first law also \emph{defines the inertial frame}: a frame of reference in which the first law holds.
Not every frame qualifies. When a bus brakes, a ball on its floor rolls forwards although nothing pushes
it; measured from the bus, the first law fails, so the braking bus is not an inertial frame. Any frame
moving at constant velocity relative to an inertial frame is also inertial. The ground is very nearly
inertial (the Earth's rotation gives accelerations of at most about \SI{0.03}{m/s^2}), and the book
treats it as inertial until \chref{ch:circdyn} studies accelerating frames.

\textbf{Inertia} is this reluctance of a body to change its velocity, and bodies have it in different
amounts: the push that sets a football rolling barely moves a car. The \textbf{mass} $m$ of a body, a
scalar measured in kilograms, is the measure of its inertia; the second law below makes the measure
precise, since the same net force gives twice the mass half the acceleration.

\begin{trapbox}{zero force means zero acceleration, not zero velocity}
$\vF_{\text{net}}=0$ gives $\va=0$: the velocity is \emph{constant}, not necessarily zero. A puck
sliding on frictionless ice at steady speed has no net force on it. ``No force'' and ``at rest'' are
different statements.
\end{trapbox}

\subsection{\texorpdfstring{The second law: $\vF=m\va$}{The second law: F = ma}}

The first law says when velocity stays constant; the second says how it changes when it does not. Newton stated it in terms of \textbf{momentum} $\vp=m\vv$ ($\vv$ the velocity), a vector along
the velocity measured in \si{kg.m/s}, which captures the everyday sense that a loaded truck at walking
pace is harder to stop than a cricket ball at the same speed.

\begin{keybox}{Newton's second law}
The rate of change of linear momentum is proportional to the net (resultant) force and is in its
direction. With units chosen so that the constant is 1 (this choice defines the newton,
$\SI{1}{N}=\SI{1}{kg.m/s^2}$), and with $\vF$ here standing for the net force,
\[
\vF=\frac{d\vp}{dt}.
\]
For constant mass this is the familiar $\vF=m\va$, with $\va$ the acceleration.
\end{keybox}

\begin{derivbox}{From momentum to $\vF=m\va$}
$\vF=\dfrac{d\vp}{dt}=\dfrac{d(m\vv)}{dt}=m\dfrac{d\vv}{dt}=m\va$ when $m$ is constant. The momentum form
is the deeper one: it carries over into relativity, where $\vF=m\va$ does not. (For a body that gains
or loses mass, such as a rocket, the law must be applied to a fixed mass that includes the ejected gas.)
\end{derivbox}

The second law is \emph{instantaneous} (the acceleration at an instant is set by the force at that same
instant), and it holds component by component: $F_x=ma_x$, so a force along one axis changes only that
component of velocity. It contains the first law's prediction ($\vF=0\Rightarrow\va=0$), but it holds
only in an inertial frame, and it is the first law that tells us which frames those are. The laws fail at
speeds approaching that of light (relativity) and on the scale of atoms (quantum mechanics), but for
everything in this book they are accurate far beyond any measurement we shall make.

The first force to which the second law is applied is the one every body near the Earth feels.

\begin{defbox}{Mass and weight}
The \textbf{weight} $\vc{W}$ of a body is the gravitational force the Earth exerts on it. Near the
Earth's surface every freely falling body has the same acceleration $\vg$ (\chref{ch:straightline}),
so by the second law $\vc{W}=m\vg$: magnitude $W=mg$, directed vertically down. Mass is a scalar in
kilograms and belongs to the body; weight is a vector in newtons and depends on where the body is. A
\SI{60}{kg} student weighs about \SI{590}{N} on Earth and about \SI{97}{N} on the Moon, and has a mass
of \SI{60}{kg} in both places.
\end{defbox}

\subsection{Impulse: force acting over time}

The momentum form of the second law pays off at once when a force is brief. In a collision or a blow
the force is large, lasts milliseconds and is rarely known in detail; such forces are called
\textbf{impulsive}. What can be known is the force's integral over time, because that integral equals
the change in momentum: integrating $\vF_{\text{net}}=d\vp/dt$ from $t_1$ to $t_2$ gives
$\int\vF_{\text{net}}\,dt=\int d\vp=\vp_2-\vp_1$. The integral is called the \textbf{impulse} $\vc{J}$ (geometrically, the area under the force--time
curve, as Figure~\ref{fig:impulse} shows),
and the average force $\vF_{\text{avg}}$ over the interval $\Delta t=t_2-t_1$ is \emph{defined} by
$\vF_{\text{avg}}\Delta t=\vc{J}$. An ordinary force such as a weight, acting for
the same millisecond, delivers a negligible impulse by comparison, which is why it may be ignored while a
bat strikes a ball or a gun fires.

\begin{keybox}{Impulse--momentum theorem}
\[
\vc{J}=\int_{t_1}^{t_2}\vF_{\text{net}}\,dt=\vF_{\text{avg}}\,\Delta t=\Delta\vp.
\]
Impulse is a vector equal to the change in momentum; its unit is $\si{N.s}=\si{kg.m/s}$.
\end{keybox}

\begin{figure}[htbp]
\centering
\tikzfig{impulse_ft}
\figcaption{fig:impulse}{Impulse is the area under the force--time curve. In a crash the change in momentum is fixed by the speed before and the zero speed after. An airbag spreads the same area over a much longer time (the figure compresses the real ratio), so $F_{\text{avg}}=\Delta p/\Delta t$ is smaller.}
\end{figure}

\why{Every crumple zone, helmet and bent knee on landing applies this one idea: for a fixed $\Delta p$, a
longer contact time $\Delta t$ means a smaller average force. You cannot cheat the momentum change, but
you can spread it out, and force, not momentum, is what breaks bones.}

\subsection{The third law}

The first two laws describe one body responding to the forces on it. They leave open where those forces
come from, and the answer is always another body; the third law says how the two bodies' forces are
related. It is the most frequently misapplied statement in mechanics, so every word of it matters.

\begin{keybox}{Newton's third law}
If body A exerts a force $\vF_{AB}$ on body B, then B exerts a force $\vF_{BA}=-\vF_{AB}$ on A. The two
forces are equal in magnitude, opposite in direction, of the same type and simultaneous.
\end{keybox}

\begin{remarkbox}{Where the third law holds}
It holds for contact forces, gravitation and the electric forces between charges at rest. For magnetic
forces between moving charges it fails as stated (the electromagnetic field then carries momentum).
Nothing in this book depends on that exception.
\end{remarkbox}

\begin{figure}[htbp]
\centering
\tikzfig{thirdlaw_book}
\figcaption{fig:thirdlaw}{The most misapplied law in mechanics. A book rests on a table. Its weight $W$ and the table's push $N$ both act \emph{on the book} (left), so they are \emph{not} a third-law pair; they are equal only because the book is in equilibrium. The genuine pairs (right, matched by colour) are $N$ with the book's push $N'$ on the table, and $W$ with the book's pull $W'$ on the Earth. Each pair acts on two \emph{different} bodies, which is why its members can never cancel.}
\end{figure}

\begin{trapbox}{action--reaction pairs never act on the same body}
The two forces of a pair (often called \emph{action} and \emph{reaction}) always act on two different
bodies, as in Figure~\ref{fig:thirdlaw}. If two forces act on the same body, they are not a third-law
pair, however equal they look.
\end{trapbox}

Third-law pairs drive nearly all propulsion. A walker's foot pushes back on the ground through
friction, and the ground pushes the walker forwards; a rocket pushes its exhaust gas downwards and the
gas pushes the rocket upwards. In each case the force that moves you is the one acting \emph{on} you,
and its partner acts on the ground or the gas.

Two further misreadings are worth naming. The law makes no exception for size: when a truck hits a
small car, each exerts a force of the same magnitude on the other, and the car fares worse only because
that force gives its smaller mass a larger acceleration. And the puzzle ``if the cart pulls back on the
horse as hard as the horse pulls the cart, how can either move?'' dissolves once each body is considered
separately, with only the forces acting on it. The cart accelerates because the horse's pull on it exceeds
the resistance at its wheels; the horse accelerates because the ground's forward friction on its hooves
exceeds the cart's backward pull on it.

\section{Momentum Conservation}

Combining the second and third laws gives the first of the conservation principles. It concerns a
\textbf{system}, any collection of bodies we choose to consider together. Forces that bodies of the
system exert on one another are \emph{internal}; forces from anything outside are \emph{external}. The
principle needs no new assumption, only the observation that internal forces come in cancelling pairs.

\begin{derivbox}{Momentum conservation for an isolated pair}
Two bodies interact with only their mutual forces; as in the third law, $\vF_{12}$ is the force body 1
exerts on body 2, and $\vF_{21}$ the force body 2 exerts on body 1. By the third law
$\vF_{12}=-\vF_{21}$ at every instant. Writing $\vF=d\vp/dt$ for each and adding:
\[
\frac{d\vp_1}{dt}+\frac{d\vp_2}{dt}=\vF_{21}+\vF_{12}=0\ \Rightarrow\ \vp_1+\vp_2=\text{constant}.
\]
Whatever momentum one body gains in any interval, the other loses.
\end{derivbox}

For a system of many bodies the same argument works pair by pair: each internal force has its
third-law partner inside the system, so the internal forces sum to zero and only external forces can
change the total momentum.

\begin{keybox}{Conservation of linear momentum}
If the net external force $\vF_{\text{ext}}$ on a system is zero, its total momentum $\vc{P}$ (the vector
sum of the momenta of its bodies) is constant:
$\vF_{\text{ext}}=0\Rightarrow\vc{P}=\sum_i m_i\vv_i=\text{constant}$. It holds component by component:
if the external force vanishes along one direction only, that component of $\vc{P}$ is still conserved.
\end{keybox}

The recoil of a gun is the cleanest case. For the gun and bullet together the horizontal external
force is negligible during the firing, and the initial momentum is zero. With subscripts $b$ for the
bullet and $g$ for the gun, $m_b\vv_b+m_g\vv_g=\vc0$: the
gun recoils opposite to the bullet, with speed $v_g=(m_b/m_g)v_b$, far smaller because the gun is far heavier.

\section{Equilibrium and Common Forces}

Momentum conservation concerned a system with zero net external force. The same condition applied to a
single particle is the simplest situation in dynamics. A particle is in \textbf{equilibrium} when the
net force on it is zero; by the second law its acceleration is then zero, so it stays at rest or moves
with constant velocity. (An extended body that can turn needs a second condition, on torques, which
\chref{ch:rotation} supplies.) This section introduces the two forces besides weight that appear in
almost every problem, then a tool for three forces in balance.

\subsection{The normal force and tension: two constraint forces}

Weight has a formula, $mg$. The two other commonest forces in mechanics do not: they take whatever value
is needed to enforce a \emph{constraint} (a surface cannot be entered, a string cannot stretch), and
finding that value is part of solving the problem.

\begin{defbox}{Normal force}
The normal force $N$ (a magnitude, in newtons) is the push a surface exerts perpendicular to itself
(``normal'' means perpendicular). It takes exactly the value needed to prevent the surfaces
interpenetrating, found by applying $\vF=m\va$ perpendicular to the surface. A surface can only push,
never pull, so $N\ge0$; if your equations ever demand $N<0$, the surfaces have separated.
\end{defbox}

\begin{figure}[htbp]
\centering
\tikzfig{normal_force}
\figcaption{fig:normal}{The normal force adjusts to the situation: $mg$ on the flat, $mg\cos\theta$ on an
incline, $m(g+a)$ in an upward-accelerating lift. It is whatever value stops the surfaces
interpenetrating, and it drops to zero the instant contact is lost.}
\end{figure}

Physically, no surface is perfectly rigid. The book compresses the table by an imperceptible amount, and
the compressed surface pushes back like a very stiff spring; because it is so stiff, a further
imperceptible compression supplies whatever extra push is needed. Because $N$ adjusts in this way, the
reflex $N=mg$ is a trap: it holds only on a horizontal surface with no other vertical forces and no
vertical acceleration. Figure~\ref{fig:normal} shows the same body receiving three different normal forces.

The \textbf{tension} $T$ in a string (a magnitude, in newtons) is the pull it exerts on whatever is attached to either end,
directed along the string. A string can pull but never push, so $T\ge0$; if the equations demand $T<0$,
the string has gone slack. A taut string is very slightly stretched and every piece of it pulls on its
neighbours, so tension exists at every point along it: the tension at a point is the force with which
the two sides of an imaginary cut there pull on each other. An \emph{ideal} string is light (its mass is
negligible) and inextensible.

\begin{derivbox}{Why an ideal string has one tension throughout}
Isolate a short segment of a massless string. The only forces along it are the tensions at its two ends;
since its mass is zero, $\sum F=ma=0$ forces those tensions to be equal. Sliding this argument along the
whole string shows that the tension is the same everywhere. For a string \emph{with} mass the argument
fails, and the tension varies: a heavy rope hanging from a ceiling is most taut at the top, which
supports all the rope below (a challenge problem in \chref{ch:connected} works out the tension along a heavy rope dragged across a table). A light, frictionless pulley (a wheel over which the string runs) changes the
direction of the string but not its tension (a pulley with mass, which must be set turning, does change it; \chref{ch:rotation} treats that case). Because the string is inextensible, its total length is fixed, and
that links the accelerations of the bodies it joins: for a single string over fixed pulleys they share
the same magnitude of acceleration (\chref{ch:connected} treats the general case).
\end{derivbox}

\subsection{Lami's theorem}

A common equilibrium problem is a body held by exactly three forces, such as a lamp hung from two
cords. Resolving into components always works, but for three forces there is a shortcut.

\begin{keybox}{Lami's theorem}
If three concurrent, coplanar forces hold a particle in equilibrium, each is proportional to the sine of
the angle between the other two:
\[
\frac{F_1}{\sin\alpha}=\frac{F_2}{\sin\beta}=\frac{F_3}{\sin\gamma},
\]
where $\alpha$ is the angle between $\vF_2$ and $\vF_3$, and so on.
\end{keybox}

\begin{derivbox}{Lami's theorem from the force triangle}
Because $\vF_1+\vF_2+\vF_3=\vc0$, the three vectors drawn head to tail close into a triangle (Figure~\ref{fig:lami}, right). Placing
$\vF_3$ head to tail after $\vF_2$ turns the angle $\alpha$ between them into an exterior angle of the
triangle, so the interior angle there, which is opposite $F_1$, is $\pi-\alpha$; and so on. The sine rule gives
$F_1/\sin(\pi-\alpha)=F_2/\sin(\pi-\beta)=F_3/\sin(\pi-\gamma)$, and $\sin(\pi-x)=\sin x$.
\end{derivbox}

\begin{figure}[htbp]
\centering
\tikzfig{lami}
\figcaption{fig:lami}{Three concurrent forces in equilibrium (left) close head-to-tail into a triangle
(right), because their resultant is zero. Applying the sine rule to that force triangle, whose
interior angles are the supplements of those between the forces, gives Lami's theorem directly.}
\end{figure}

\begin{exbox}{A lamp on two cords}\label{ex:lami}
A lamp of weight \SI{60}{N} hangs from two cords making $30^\circ$ and $60^\circ$ with the horizontal
ceiling. Find the tensions $T_1$ (shallower cord) and $T_2$ (steeper cord).

\stepn{1} \textbf{Angles between the forces.} The cords are $180^\circ-30^\circ-60^\circ=90^\circ$ apart.
The shallower cord is $60^\circ$ from the upward vertical, so it makes $120^\circ$ with the weight; the
steeper cord makes $150^\circ$ with it.

\stepn{2} \textbf{Lami.} $T_1/\sin150^\circ=T_2/\sin120^\circ=60/\sin90^\circ$, so $T_1=\SI{30}{N}$ and
$T_2\approx\SI{52}{N}$.

\stepn{3} \textbf{Check.} Horizontally, $30\cos30^\circ\approx26.0=52\cos60^\circ$. The steeper cord takes
more of the load; as both cords flatten towards the horizontal, both tensions grow without limit, which
is why a loaded washing line can never be pulled straight. \ans{\SI{30}{N}; \SI{52}{N}}
\end{exbox}

\section{The Free-Body Method}

The laws say what forces do, and the previous section catalogued the forces that occur most often. What
remains is a reliable way to bring the two together on a particular body. Every problem in this Part is
solved by the same procedure. It is set out here as five explicit steps, because nearly every error in dynamics
comes from skipping one of them.

\begin{methodbox}{The five-step recipe}
\textbf{1.} Isolate one body; draw only the forces acting \emph{on} it (weight, normal, tension,
friction, applied). \textbf{2.} Choose axes, along and perpendicular to the motion (along the incline
for slopes). \textbf{3.} Resolve each force and write $\sum F=ma$ per axis (set $a=0$ in equilibrium).
\textbf{4.} Add constraints (from the string length; for one string over fixed pulleys the bodies
share $|a|$). \textbf{5.} Solve, then check signs and
limits.
\end{methodbox}

\begin{trapbox}{draw only forces that have a source}
Every arrow on a free-body diagram must be exerted by an identifiable body: the Earth (weight), a
surface (normal force, friction), a string (tension), a spring. Three popular arrows have no source and
must be left out: a ``force of motion'' carrying a moving body forwards (the first law says none is
needed); $m\va$, which is the \emph{result} of the forces, not one of them; and a ``centripetal force'',
which is a role played by real forces (\chref{ch:circdyn}). Forces the body exerts \emph{on} other
things belong on their diagrams, not on its own.
\end{trapbox}

\begin{figure}[htbp]
\centering
\tikzfig{fbd_incline}
\figcaption{fig:fbd}{A free-body diagram for a block at rest on a rough incline: weight $mg$ (down, from the centre of mass), normal force $N$ (perpendicular to the surface) and friction $f$ (along it, up the slope). Resolving $mg$ along and perpendicular to the incline (dashed) turns the vector equation into two scalar ones: $f=mg\sin\theta$ and $N=mg\cos\theta$.}
\end{figure}

Figure~\ref{fig:fbd} shows the steps on a block held at rest by friction (a force treated fully in
\chref{ch:friction}); the example below runs all five on a smooth slope, where the block accelerates.

\begin{exbox}{The five steps on a smooth incline}\label{ex:fbdincline}
A \SI{2}{kg} block on a smooth incline of angle $30^\circ$ is pulled up the slope by a light string
parallel to the surface, with tension \SI{16}{N} ($g=\SI{10}{m/s^2}$). Find its acceleration and the
normal force on it.

\stepn{1} \textbf{Isolate and draw.} Three forces act on the block: its weight $mg=\SI{20}{N}$ vertically
down, the normal force $N$ perpendicular to the incline, and the tension $T=\SI{16}{N}$ up the slope.
There is no friction, because the incline is smooth.

\stepn{2} \textbf{Axes.} Take $x$ up the slope and $y$ perpendicular to it, so that the acceleration lies
along $x$ alone.

\stepn{3} \textbf{Resolve and apply $\sum F=ma$.} The weight has components $-mg\sin30^\circ=\SI{-10}{N}$
along $x$ and $-mg\cos30^\circ\approx\SI{-17.3}{N}$ along $y$. Along $x$: $16-10=2a$, so
$a=\SI{3}{m/s^2}$ up the slope. Along $y$: $N-17.3=0$, so $N\approx\SI{17.3}{N}$.

\stepn{4} \textbf{Constraints.} The block stays on the surface, which is why $a_y=0$; there is no second
body.

\stepn{5} \textbf{Check.} With $T=0$ the $x$ equation gives $a=-g\sin30^\circ=\SI{-5}{m/s^2}$, the familiar
slide down a smooth incline; and $N<mg$, as it must be on a slope. \ans{\SI{3}{m/s^2}; \SI{17.3}{N}}
\end{exbox}

\subsection{Apparent weight: what a weighing scale actually measures}

The same method settles the third misconception named at the start of the chapter.
A weighing scale does not measure gravity. It measures the \textbf{normal force} it must exert to
support you, and that equals $mg$ only when you are not accelerating vertically. A lift is the cleanest
demonstration, and Figure~\ref{fig:lift} sets out its four cases.

\begin{figure}[htbp]
\centering
\tikzfig{lift_apparent}
\figcaption{fig:lift}{The four lift cases. The scale reads $N$. Accelerating up requires $N>mg$;
accelerating down requires $N<mg$; in free fall $N$ vanishes and the person is weightless.}
\end{figure}

\begin{derivbox}{Apparent weight in an accelerating lift}
Two vertical forces act on a person of mass $m$: the normal force $N$ from the floor (up) and the weight
$mg$ (down). Taking \emph{up} as positive, with $a$ the lift's (and the person's) acceleration, Newton's
second law gives
$N-mg=ma$, so
\[
\boxed{\;N=m(g+a)\;}
\]
with $a$ carrying its own sign. Reading off the four cases:
\begin{enumerate}
\item \textbf{At rest or constant velocity} ($a=0$): $N=mg$, the true weight. A smooth lift at
      cruising speed feels exactly like standing still.
\item \textbf{Accelerating upwards} ($a>0$): $N=m(g+a)>mg$; you feel \emph{heavier}. This is the
      pressed-into-the-floor sensation as a lift starts up (or as a descending lift slows to a stop).
\item \textbf{Accelerating downwards} ($a<0$): $N<mg$; you feel \emph{lighter}, the stomach lurch of a
      lift starting down (or of a rising lift slowing to a stop).
\item \textbf{Free fall} ($a=-g$): $N=0$. \textbf{Weightlessness}, not because gravity has switched
      off, but because nothing is pushing back on you. This is the condition of an astronaut in orbit,
      who is in permanent free fall around the Earth (\chref{ch:gravitation}).
\end{enumerate}
\end{derivbox}

\begin{exbox}{Reading the scale}
A \SI{50}{kg} person stands on a scale in a lift ($g=\SI{10}{m/s^2}$). What does it read when the lift
accelerates (a) upwards at \SI{2}{m/s^2}, (b) downwards at \SI{2}{m/s^2}?

\stepn{1} \textbf{(a)} $N=m(g+a)=50(12)=\SI{600}{N}$; a scale calibrated in kilograms shows
\SI{60}{kg}.

\stepn{2} \textbf{(b)} Now $a=\SI{-2}{m/s^2}$, so $N=m(g+a)=50(8)=\SI{400}{N}$, which the scale shows
as \SI{40}{kg}.

\stepn{3} \textbf{Check the limit:} at $a=g=\SI{10}{m/s^2}$ downwards the formula gives $N=0$, free
fall, as it must. \ans{\SI{600}{N}; \SI{400}{N}}
\end{exbox}

\section{Summary}\label{sec:summary7}
The first law defines force, and the inertial frames in which the laws hold. The second law,
$\vF=d\vp/dt$, reduces to $\vF=m\va$ for constant mass and is applied component by component. The third
law pairs every force with an equal and opposite force on a \emph{different} body, so the two never
cancel each other.

\begin{itemize}
\item Mass (kg, scalar) measures inertia; weight $W=mg$ (N, vector) is the Earth's pull and changes from place to place.
\item Zero net force means zero acceleration, not zero velocity.
\item Impulse $\vc{J}=\int\vF\,dt=\Delta\vp$: a force acting briefly changes momentum by its time integral.
\item With no net external force, total momentum is conserved, component by component.
\item The normal force and tension are constraint forces: they have no formula of their own and take
      whatever value $\vF=m\va$ requires. A surface can push but not pull; a string can pull but not push.
\item A scale reads the normal force, not the weight: $N=m(g+a)$ in a lift with upward acceleration $a$.
\item Lami's theorem: three concurrent forces in equilibrium satisfy $F_1/\sin\alpha=F_2/\sin\beta=F_3/\sin\gamma$.
\item The free-body method: isolate, draw real forces only, choose axes, resolve, apply $\sum F=ma$ per axis.
\end{itemize}

\section{Problem Types}\label{sec:types7}
Six problem types cover the Newton's-law questions that examinations set most often, from the general
free-body problem to settings that recur year after year. All six reward one habit above any formula:
draw the free-body diagram for each body separately before writing a single equation.

\archq{Newton's-law problems}
{Forces and masses given; find acceleration, tension, or a normal force.}
{Draw a free-body diagram for each body. Resolve along and perpendicular to the motion. Write
$\sum F=ma$ per axis; add the connecting constraint; solve simultaneously.}
{Forgetting that $N$ is not always $mg$, and treating a third-law pair as if it acted on one body. Both
are diagram errors, caught by drawing forces on \emph{one} body at a time.}
{A \SI{2}{kg} block on a frictionless table is pulled by a string with tension \SI{10}{N} directed at
$60^\circ$ above the horizontal ($g=\SI{10}{m/s^2}$). Its acceleration is:}
{\SI{5}{m/s^2}}{\SI{4.3}{m/s^2}}{\SI{20}{m/s^2}}{\SI{2.5}{m/s^2}\archans{D}{Only the horizontal
component accelerates the block: $a=F\cos60^\circ/m=10(0.5)/2=\SI{2.5}{m/s^2}$. The vertical component
$F\sin60^\circ\approx\SI{8.7}{N}$ only reduces the normal force, to $N\approx\SI{11.3}{N}$ (not $mg$), and
does not lift the block. Option (A) ignores the angle, (B) uses $\sin$ in place of $\cos$, and (C)
multiplies by the mass instead of dividing.}}

\archq{Apparent weight in a lift}
{``A person of mass $m$ stands on a scale in a lift accelerating up or down; what does the scale
read?''}
{The scale reads the \emph{normal force}, not $mg$. Write $N-mg=ma$ with \emph{up} positive, so
$N=m(g+a)$ with $a$ carrying its own sign: $a>0$ going up gives $N>mg$; $a<0$ gives $N<mg$; free fall
$a=-g$ gives $N=0$.}
{Assuming a lift moving \emph{upwards} always means a larger reading. It is the \emph{acceleration}
that matters, not the velocity: a lift rising at constant speed reads exactly $mg$.}
{A \SI{40}{kg} boy stands on a scale in a lift accelerating downwards at \SI{2.5}{m/s^2}
($g=\SI{10}{m/s^2}$). The reading is:}
{\SI{100}{N}}{\SI{300}{N}}{\SI{400}{N}}{\SI{500}{N}\archans{B}{$N=m(g-a)=40(10-2.5)=\SI{300}{N}$;
he feels lighter. Option (D), \SI{500}{N}, would be the upward-accelerating case, (C) ignores the
acceleration, and (A) is $ma$ alone, forgetting gravity.}}

\archq{Impulse and change of momentum}
{``A ball of mass $m$ strikes a wall or bat at speed $u$ and rebounds at $v$; find the impulse, or the
average force given the contact time.''}
{Impulse $=\Delta\vp=m\vv_f-m\vv_i$, and $F_{\text{avg}}=J/\Delta t$. Because the ball \emph{reverses}, put in
the signs: for a normal rebound $J=m(v+u)$, not $m(v-u)$.}
{Subtracting the speeds instead of adding them. A reversal means the momentum change is the
\emph{sum} of the magnitudes.}
{A \SI{0.5}{kg} ball hits a wall normally at \SI{10}{m/s} and rebounds at \SI{10}{m/s}. The magnitude of
the impulse on the wall is:}
{zero}{\SI{5}{N.s}}{\SI{10}{N.s}}{\SI{20}{N.s}\archans{C}{Taking towards the wall as positive,
$\Delta p=(-10)(0.5)-(+10)(0.5)=\SI{-10}{kg.m/s}$, magnitude \SI{10}{N.s}, and by the third law the wall receives an equal and
opposite impulse. Option (A) is the trap of thinking equal speeds mean no change.}}

\archq{Two blocks pushed together, the contact force}
{``Two blocks in contact on a smooth floor are pushed by a force $F$; find the acceleration and the
force between them.''}
{\textbf{Two steps.} First treat the pair as one body (the forces between the blocks are an internal
third-law pair and cancel) to get $a=\dfrac{F}{m_1+m_2}$. Then isolate the
\emph{front} block (the one not being pushed directly): the only horizontal force on it is the contact
force, so $N=m_2a$.}
{Taking the contact force to equal $F$. It is always \emph{less}, because part of $F$ goes into
accelerating the block that is pushed.}
{Blocks of \SI{2}{kg} and \SI{3}{kg} in contact on a smooth floor are pushed by \SI{10}{N} applied to the
\SI{2}{kg} block. The contact force between them is:}
{\SI{4}{N}}{\SI{6}{N}}{\SI{10}{N}}{\SI{2}{N}\archans{B}{$a=\dfrac{10}{5}=\SI{2}{m/s^2}$; then for the
\SI{3}{kg} block alone, $N=3(2)=\SI{6}{N}$. Option (A) is the force needed for the \SI{2}{kg} block, and
(C) is the trap above.}}

\archq{Recoil of a gun / explosion}
{``A gun fires a bullet; find the recoil speed,'' or ``a body at rest explodes into fragments.''}
{Momentum conservation with an initial total of zero: $m_1v_1+m_2v_2=0$, so the two move
\emph{oppositely} with $v_g=\dfrac{m_bv_b}{m_g}$. For an explosion, all fragment momenta must sum
to $\vc0$.}
{Assuming the energy of motion (kinetic energy, \chref{ch:energy}) is conserved. The burning propellant
\emph{adds} energy, so it is not; only momentum is, because the forces of the explosion are internal.}
{A \SI{10}{g} bullet is fired at \SI{500}{m/s} from a \SI{5}{kg} gun. The recoil speed of the gun is:}
{\SI{1}{m/s}}{\SI{0.5}{m/s}}{\SI{2}{m/s}}{\SI{5}{m/s}\archans{A}{$m_gv_g=m_bv_b$, so
$v_g=\dfrac{0.010(500)}{5}=\SI{1}{m/s}$, directed backwards. Option (D) is the bullet's momentum,
\SI{5}{kg.m/s}, mistaken for the gun's speed; (B) and (C) are slips of a factor of two.}}

\archq{A rope, chain or spring balance: where is the tension?}
{``Find the tension at a point in a hanging chain,'' or ``what does a spring balance read between two
masses?''}
{A spring balance reads the \emph{tension} in the string through it, not the weight hung on it. For a
hanging system with upward acceleration $a$, isolate the portion \emph{below} the cut:
$T=m_{\text{below}}(g+a)$.}
{Assuming the balance reads $mg$ regardless. In an accelerating lift, or with masses on both sides, it
does not.}
{A monkey of mass \SI{20}{kg} climbs \emph{up} a light rope with an upward acceleration of
\SI{2}{m/s^2} ($g=\SI{10}{m/s^2}$). The tension in the rope is:}
{\SI{160}{N}}{\SI{200}{N}}{\SI{240}{N}}{\SI{40}{N}\archans{C}{Climbing \emph{up} with acceleration $a$
requires $T-mg=ma$, so $T=m(g+a)=20(12)=\SI{240}{N}$. Option (A), \SI{160}{N}, is the tension if the
monkey slid down with a downward acceleration of \SI{2}{m/s^2}, (B) holds only at constant
velocity, and (D) is $ma$ alone. The monkey accelerates up because it pulls down on the rope, so the
rope pulls up on it with $T$.}}
